
\textbf{Reward models~(RMs) and verifiers.} Conventionally, RMs and verifiers are trained as discriminative models via binary classification: given a prompt and a corresponding solution or a pair of solutions), the model is either trained to predict the correctness of the solution~\citep{cobbe2021training, lightman2023let, wang2023math, uesato2022solving, luo2024improve, yu2024ovm} or a preference between the two solutions~\citep{stiennon2020learning, nakano2021webgpt}. Concretely, the RM directly produces a numerical continuous-valued score, which is then plugged into a classification objective~\eqref{eq:disc}. As such, discriminative verifiers do not utilize the generation capabilities of LLMs. In contrast to discriminative RMs, \genV\ represents the correctness decision using the log probability of specific tokens, for example `Yes' and `No'. Posing verification as generating ``yet another token'' allows it to tap better into the generation capabilities of LLMs, by making it straightforward to employ CoT reasoning and additional inference-time compute for better verification. 

\textbf{LLM-as-a-Judge.} Another line of work that poses verification as next-token prediction simply \emph{prompts} off-the-shelf LLMs to act as a verifier when provided with a rubric and a template for grading~\citep{zheng2024judging, bai2022constitutional, kim2023prometheus, deductive-verification} or many-shot ICL examples~\citep{agarwal2024many}, but \emph{without} any specific training for the same. Perhaps unsurprisingly, we find in our experiments that using more powerful LLMs (Gemini 1.0 Pro) as a judge is worse than our trained \genV\ using weaker Gemma models~(\autoref{fig:overall-result}, \ref{fig:inference_compute_scaling}), highlighting the necessity of \emph{training} generative verifiers. Our generative verifiers also exhibit good out-of-distribution generalization, which might be due to better calibrated uncertainty estimates from training~\citep{kapoor2024large}. More generally, even the strong proprietary LLMs, such as GPT-4~\citep{gpt-4-technical-report} and Gemini ~\citep{team2024gemini}, fall behind trained RMs on popular leaderboards~\citep{rewardbench}, and this gap is much larger for reasoning. 

\textbf{Using CoTs for reward models.} Prior works have also used critiques or CoT to extract preference and verification signals using LLM-as-a-Judge~\citep{self-reward-lm, meta-rewarding-lm, self-taught-evaluator}; in contrast to these works, \genV\ utilizes model-generated CoTs directly for training the verifier. Upon inference, a \genV-CoT produces its own CoTs, which it then uses to make decisions on correctness, unlike \citet{improve-rm-synthetic-critique} that simply uses CoTs from a separate highly-capable LLM. In contrast to prior work that utilizes high-quality data from humans to train critique models~\citep{self-critique-models} or train \emph{discriminative} RMs for generating code critiques~\citep{critic-gpt}, we show that \genV\ can be trained from purely synthetic, model-generated critiques. 
Concurrent work \citep{ankner2024critique} trains an RM to produce response critiques for preference pairs generated using a much more capable LLM, which are then passed as input into a RM head, separate from the base LLM. Unlike \genV\ which uses next-token prediction, their RM head is trained discriminatively akin to standard RMs. While this approach allows them to leverage CoT, it does \emph{not} allow them to unify solution generation and verification as a result of a discriminative RM head, which \genV\ seamlessly enables (Section~\ref{sec:unification}). Moreover, their synthetic critiques are not filtered for correctness, which would lead to poor verification CoTs on reasoning tasks~(\S\ref{sec:cot-verifier}). 


\textbf{Unified generation and verification.} One of the hallmark properties of \genV\ is that the same generative verifier can be co-trained with a generation objective~\eqref{eq:co-train}: when given a problem, the model is trained to produce a solution, whereas when given a problem and a candidate solution, it is trained to verify this candidate. This is related to DPO~\citep{rafailov2024direct} and its application to learning verifiers in reasoning~\citep{hosseini2024v}, which aims to unify generation (policy) and verification (reward models) by representing the reward implicitly using the logits of a policy and training the policy with a reward-modeling loss. 
For reasoning, this type of model tying has been shown to exhibit erroneous extrapolation and degradation in learned representations, which prior work has attempted to address with additional techniques~\citep{iterative-reasoning-preference-op, rl-on-incorrect,pal2024smaug,yang2024regularizing}. Of these, while \citet{yang2024regularizing} train a reward model with an auxiliary generative SFT loss, note that this loss is applied on a separate head for regularization purposes and is discarded after training; unlike \genV\ no text is produced when querying the RM. In addition, compared to DPO, {\genV} uses a simpler next-token prediction loss, does not require a reference policy, and obtains significantly better verification performance~(\autoref{fig:overall-result}, \ref{fig:baselines}).
